\begin{table}[t]
\centering
\caption{\textbf{Impact of Auxiliary Tasks on Generation Performance.} Comparison of MAE and Jaccard Index w/ and w/o edge reconstruction and node classification tasks in the GNN. Including both tasks yields the best alignment with the scene graph.}
\vspace{-1em}
\scriptsize 
\setlength{\tabcolsep}{9pt}
\begin{tabular}{cc|cc}
\toprule
Reconstruction            & Classification            & MAE (\(\downarrow)\)           & Jaccard       \\ \midrule \midrule
\ding{51} & \ding{51} & \textbf{0.63} & \textbf{0.93} \\
\ding{51} & \ding{55}     & 0.89          & 0.84          \\
\ding{55}     & \ding{51} & 0.80          & 0.83          \\
\ding{55}     & \ding{55}     & 0.79          & 0.81         \\ \bottomrule
\end{tabular}
\label{tab:different_tasks}
\vspace{5mm}
\end{table}

\begin{table}[t]
\centering
\caption{\textbf{Comparison of Training Strategies for Our Method.} The bolded row is our adopted strategy.}
\vspace{-1em}
\scriptsize 
\setlength{\tabcolsep}{9pt}
\resizebox{0.47\textwidth}{!}{
\begin{tabular}{ccc|cc}
\toprule
Diffusion        & GNN          & LOC             & MAE (\(\downarrow)\) & Jaccard       \\ \midrule \midrule
Pre-train        & Pre-train        & Pre-train           & 0.79                 & 0.88          \\
Scratch          & Scratch          & Scratch             & 1.01                 & 0.82          \\
Scratch          & Pre-train        & Pre-train           & 0.95                 & 0.90          \\
\textbf{Scratch} & \textbf{Scratch} & \textbf{Post-train} & \textbf{0.63}        & \textbf{0.93} \\ \bottomrule
\end{tabular}}
\label{tab:different_stratigies}
% \vspace{-2em}
\end{table}

\section{Conclusion}
In this work, we propose a solution that integrates an interactive system, BEV Embedding Map, and diffusion generation to enable controllable 3D outdoor scene generation. The challenges stem from complex outdoor landscapes with rich information and structural diversity. Our approach utilizes scene graphs to transition information from sparse to dense representations. Coupled with the interactive system, it enables users to intuitively and concisely generate their desired 3D outdoor scenes. Comparative experiments demonstrate that our method achieves more accurate object quantities and alignment with the input scene graph. These results indicate that our approach is a robust and effective solution for controllable 3D outdoor scene generation.
\label{sec:conclusion}